% !PRINTED 261
\weberpage{261}

III.\ \gesperrt{If $R$ is different from zero, one can arrange the indices $1, 2 \ldots m$ so that in the chain of principal subdeterminants}
\weq{3}{R_m,\ R_{m-1} \ldots R_1,\ R_0}
\gesperrt{no two consecutive terms vanish.}

The determinant relation (2) yields, if one sets
\[
A = R_{k+1}, \quad (k = 1, 2 \ldots m-1)
\]
an equation between three consecutive terms of the chain (3), which we can write thus:
\weq{4}{R_k S_k - T_k^2 = R_{k-1} R_{k+1},}
in which $S_k$ and $T_k$ are certain subdeterminants of $R$, and hence integral rational functions of the coefficients $a_{i,k}$.

More precisely described, $R_k$, $S_k$, $T_k$ are first subdeterminants of $R_{k+1}$, and indeed
\[
R_k = \frac{\partial R_{k+1}}{\partial a_{k+1,k+1}}, \quad S_k = \frac{\partial R_{k+1}}{\partial a_{k,k}}, \quad T_k = \frac{\partial R_{k+1}}{\partial a_{k,k+1}}.
\]

An arrangement of the indices corresponding to theorem (3) we shall suppose chosen for what follows. Then, if $R_k$ vanishes, $R_{k-1}$ and $R_{k+1}$ are different from zero, and (4) shows that they have opposite signs; hence:

IV.\ \gesperrt{If an inner term of a chain of principal subdeterminants vanishes, then the two adjoining terms have opposite signs.}

\begin{center}
\S.\,83.\\[4pt]
\gesperrt{Number of the positive and negative squares.}
\end{center}

In order to determine the number of positive and negative squares of a form with non-vanishing determinant, we assume first that in the chain of principal subdeterminants
\[
R_m,\ R_{m-1},\ R_{m-2} \ldots R_1,\ R_0 = 1
\]
no term vanishes. We can then determine the coefficient $\lambda$ so that the determinant of the form
%CONT>>
